% =====================================================================
%  ANEXO C: INSTRUMENTO DE RECOLECCIÓN
% =====================================================================
\anexo{C}{Instrumento de Recolección}

El formulario tiene tres partes: el consentimiento informado, una verificación de que la persona es mayor de edad y está matriculada, y un campo para escribir el relato. No pide nombres, capturas de chats, números de contacto, códigos de estudiante ni datos de terceros.

\vspace{0.4cm}
\noindent\textbf{Consigna:}

\begin{tcolorbox}[colback=gray!5, colframe=gray!50, arc=2pt, boxrule=0.5pt]
\itshape
Si deseas participar, cuéntanos con tus propias palabras una situación de tu vida universitaria que quieras compartir. Puede ser algo que te haya resultado incómodo o problemático, o simplemente una interacción cotidiana. Explica qué pasó y el contexto que creas necesario, sin incluir nombres, fechas exactas ni detalles que permitan reconocer a alguien. Puedes dejar de responder en cualquier momento.
\end{tcolorbox}

\vspace{0.4cm}
Al final del formulario se indica que este no es un canal de denuncia y se incluyen los datos de contacto de las oficinas de orientación y atención de la universidad.
